\thispagestyle {empty}

\vspace{1cm}
\begin{center}
    \small
    \onehalfspacing
    { 
    МИНОБРНАУКИ РОССИИ

    Федеральное государственное бюджетное образовательное\\учреждение высшего профессионального образования

    <<САНКТ-ПЕТЕРБУРГСКИЙ ГОСУДАРСТВЕННЫЙ\\МОРСКОЙ ТЕХНИЧЕСКИЙ УНИВЕРСИТЕТ>>
    }
\end{center}

\vspace{2cm}

\begin{center}
    \onehalfspacing
    { 
    Н.Ю. Часовников
    }
\end{center}

\vspace{2cm}

\begin{center}
    \huge
    \onehalfspacing
    { 
    РАСЧЕТЫ ПО СТАТИКЕ
    
    МОРСКИХ СУДОВ
    }
\end{center}

\vspace{1cm}

\begin{center}
    \onehalfspacing
    { 
    Учебное пособие
    }
\end{center}

\vspace*{3cm}

\begin{center}
    \onehalfspacing
    { 
    Санкт-Петербург

    2026
    }
\end{center}

\newpage

\thispagestyle {empty}
УДК 629.12

ББК 39.42-011

\begin{center}
    Рецензенты:

    \vspace{0.5cm}

    \textit{доктор технический наук, доцент} \textbf{Ярисов~В.В.}

    \textit{кандидат технических наук} \textbf{Ишков~В.В.}

\end{center}

\vspace{1.0cm}

\textbf{Часовников~Н.Ю.}

Расчеты по статике морских судов: учебное пособие. --- СПбГМТУ. --- СПб., 2026. --- 126~с.

\vspace{0.5cm}

\textbf{ISBN}

\vspace{0.5cm}

В настоящем учебном пособии изложены методические указания по выполнению основных видов расчетов по статике морских судов.
Все расчетные схемы, предназначенные для реализации в табличных редакторах, снабжены поясняющим графическим материалом.

Учебное пособие предназначено для студентов факультета кораблестроения и океанотехники, обучающихся по очной, вечерней и заочной формам обучения, и выполняющих домашние работы и курсовые проекты по курсам <<Статика корабля>>, <<Теория корабля>>, <<Теория и устройство корабля>>.

\vspace{3.5cm}

{\footnotesize
\copyright СПбГМТУ, 2026}
